\section{Luke 24:12 Beside John 20:3--10: Textual Status and Harmony}
\label{sec:luke-john}

When St. Luke's Gospel narrates the morning of the Resurrection, it devotes only a single verse to Peter's visit to the tomb---yet that single verse carries both a celebrated textual history and a real narrative puzzle beside John 20:3--10.

In the Greek text of Luke 24:12 as printed in the Byzantine tradition (Robinson--Pierpont), the sentence runs:
\begin{quote}
\latin{Ho de Petros anastas edramen epi to mn\={e}meion, kai parakypsas blepei ta othonia keimena mona: kai ap\={e}lthen pros heauton thaumaz\={o}n to gegonos.}
\end{quote}
The Douay--Rheims renders the corresponding Clementine Vulgate (\latin{Petrus autem surgens cucurrit ad monumentum: et procumbens vidit linteamina sola posita, et abiit secum mirans quod factum fuerat}): ``But Peter rising up, ran to the sepulchre and, stooping down, he saw the linen cloths laid by themselves: and went away wondering in himself at that which was come to pass.''

\subsection{A Western Non-Interpolation}

Before comparing Luke 24:12 with John 20:3--10, textual candor requires noting the manuscript history of Luke 24:12 itself. In 1882, Brooke Foss Westcott and Fenton John Anthony Hort placed Luke 24:12 inside double brackets in their Greek New Testament as one of their famous ``Western non-interpolations''---a small group of passages near the end of St. Luke's Gospel that are present in Alexandrian and Byzantine manuscripts but absent from Codex Bezae (D) and several Old Latin manuscripts. In their Appendix (\emph{Notes on Select Readings}, p. 71), Westcott and Hort summarize the external testimony and their own textual judgment plainly:
\begin{quote}
``xxiv 12 [\latin{Ho de Petros \dots\ gegonos.}] $<$ Western (Gr. Lat.), D a b e rhe Eus.Can; not c ff syr.vt Eus.Mar (distinctly). Omitted likewise at the beginning of one lection (first hand) in syr.hr, and in the harmonistic narrative of fu; but probably in both cases by accident. Text from Jo xx 3--10 (except \latin{anastas} and \latin{thaumaz\={o}n to gegonos}), condensed and simplified, with omission of all that relates to the other disciple. A Western non-interpolation.''
\end{quote}

Even Westcott and Hort's own apparatus shows that the verse is far from a late Byzantine addition: it is absent from D, Old Latin manuscripts \latin{a, b, e, r\textsuperscript{1}}, and Eusebius's Canon Tables, yet attested by other Old Latin manuscripts (\latin{c, ff\textsuperscript{2}}), the Old Syriac (\latin{syr.vt}), and Eusebius's \emph{Ad Marinum} (``distinctly''), as well as the great Alexandrian uncials and the Latin Vulgate.\endnote{In twentieth-century textual criticism, the discovery of Papyrus 75 (Bodmer XIV--XV, early third century), which contains Luke 24:12, decisively strengthened the case for the verse's authenticity alongside Codex Vaticanus (B) and Codex Sinaiticus; accordingly, modern critical editions such as the SBL Greek New Testament (Holmes v1.2) print Luke 24:12 in the text inside half-brackets rather than excluding it. For Catholic readers, the verse stands in the canonical Vulgate text received by the Church.}

Two smaller internal variants within Luke 24:12 also matter for careful reading:

\begin{itemize}
\item \textbf{With or without \latin{keimena} (``lying''):} The Byzantine textform reads \latin{ta othonia keimena mona} (``the linen cloths lying alone,'' matching the Clementine Vulgate's \latin{linteamina sola posita} and Douay--Rheims's ``the linen cloths laid by themselves''). The SBL Greek New Testament reads simply \latin{ta othonia mona} (``the linen cloths alone,'' without \latin{keimena}), which the Revised Version of 1895 reflects (``he seeth the linen cloths by themselves'').
\item \textbf{Where \latin{pros hauton} attaches:} The closing clause reads \latin{kai ap\={e}lthen pros hauton thaumaz\={o}n to gegonos} (Byzantine \latin{pros heauton}). Taken with the participle \latin{thaumaz\={o}n}, it yields the Vulgate's \latin{abiit secum mirans} and Douay--Rheims's ``went away wondering in himself.'' Taken with the main verb \latin{ap\={e}lthen}, it matches the exact idiom of John 20:10 (\latin{ap\={e}lthon \dots\ pros hautous}, ``departed to their home'') and yields the Revised Version's ``he departed to his home, wondering at that which was come to pass.'' Both readings are grammatically possible in Koine Greek, and both leave Peter in the same interior posture: marveling at what had happened, not yet proclaiming the Resurrection.
\end{itemize}

\subsection{Comparing the Two Accounts: St. Augustine and St. Bede}

Read side by side, Luke 24:12 and John 20:3--10 share striking verbal links---the verb \latin{parakypsas} (``stooping down''), the present \latin{blepei} (``sees''), the exact noun \latin{ta othonia} (``the linen cloths''), and the departure \latin{pros hauton} / \latin{pros hautous}---while differing on three visible points:

\begin{enumerate}
\item Luke names only Peter running to the tomb in verse 12, whereas John narrates Peter and the beloved disciple running together.
\item Luke says that Peter stooped down (\latin{parakypsas}) and saw the linen cloths alone, without mentioning his entrance or the separate \latin{soudarion}; John assigns the stooping (\latin{parakypsas}) to the beloved disciple, and has Peter enter the tomb and see both the \latin{othonia} and the \latin{soudarion}.
\item Luke places Peter's run after the women report the angelic message (Luke 24:9--11), whereas John places Peter and the beloved disciple's run immediately after Mary Magdalen's first report of the removed stone, before her own vision of the angels and the risen Lord (John 20:1--18).
\end{enumerate}

The Church's Latin exegetes did not gloss over these divergences; they examined them sentence by sentence. First, St. Augustine pointed out in \emph{The Harmony of the Gospels} (III.25.70) that Luke's own chapter confirms John's report of more than one disciple visiting the tomb. Later that afternoon on the road to Emmaus, Cleophas and his companion tell the unrecognized stranger: ``And some of our people went to the sepulchre and found it so as the women had said: but him they found not'' (Luke 24:24, \latin{tines t\={o}n syn h\={e}min}, ``some of those with us''---plural). Luke singled out Peter in verse 12 by prominence as the chief Apostle, Augustine explains, just as John singled out Mary Magdalen at the head of the women while preserving her plural ``we know not where they have laid him'' (John 20:2).

Second, Augustine directly resolved the sequence of stooping and entering (\emph{Harmony} III.25.70):
\begin{quote}
``A difficulty, however, may also be felt in the circumstance that the same Luke does not say that Peter entered, but only that he stooped down and saw the linen clothes laid by themselves, and that thereupon he departed, wondering in himself; whereas John intimates that it was rather himself (for he is the disciple whom Jesus loved) that looked at the scene in this fashion, not going within the sepulchre, which he was the first to reach, but simply bending down and beholding the linen clothes laid in their place; although he also adds that he did enter the tomb afterwards. The explanation, therefore, is simply this, that Peter at first did stoop down and look in after the fashion which Luke specifies, but to which John makes no allusion; and that he went actually in somewhat later, but still before John entered.''
\end{quote}

Three centuries later, St. Bede the Venerable took up the exact same question in his \emph{Exposition of the Gospel of Luke} (Book VI, on Luke 24:12), repeating Augustine's insight with brisk clarity: ``Luke briefly mentions Peter's run, while John sets out more fully how this was done, saying that the disciple whom Jesus loved also ran with Peter, indicating himself. \dots\ It must be understood that Peter first stooped down and saw what Luke mentions and John passes over, and afterward went in so that he might discern the interior more carefully---yet entering before John went in.''\endnote{St. Bede the Venerable, \emph{In Lucae evangelium expositio}, Book VI, on Luke 24:12 (Migne, PL 92, col. 625B): \latin{``Sed intelligendum est Petrum primo procumbentem vidisse quod Lucas commemorat, Joannes tacet; post autem ingressum, ut diligentius interiora dignosceret, ingressum, tamen antequam Joannes intraret.''}}

As for timing, Augustine treats Luke 24:12 as a brief recapitulation inserted after the women's report, whereas John preserves the step-by-step morning sequence. Whatever exact choreography one adopts for the women's visits at dawn, one theological point stands out with equal force in both Gospels: Luke 24:12 says nothing about a folded table napkin, mentions no \latin{soudarion} at all, and leaves Simon Peter departing from the empty tomb in silent amazement (\latin{thaumaz\={o}n to gegonos})---a wonder that will reach full apostolic proclamation only when the risen Lord personally appears to Simon later that day (Luke 24:34; 1 Cor 15:5).
